\clearpage
\section{Especificação formal M10: patrimônio}
\label{sec:formal-m10}
\sloppy

M10 formaliza o alvo normativo de patrimônio: a identidade entre o modelo
catalográfico (\texttt{Resumo\_Bem\_Patrimonial}) e a unidade física
(\texttt{Bem\_Patrimonial}), a plaqueta canônica, a unicidade permanente, a
máquina de estados, o versionamento, as requisições de adição/edição, os locks,
o histórico, a baixa e a composição abstrata com M7 (idempotência) e M9
(autorização). A evidência cobre somente o modelo formal declarado.

\subsection*{CUE}
O contrato CUE (\texttt{\#M10Contrato}) verifica shapes, enums fechados,
obrigatoriedade, nulabilidade e limites locais de bem, resumo, requisição de
adição, requisição de edição, lock, baixa, evento histórico e alteração. O CUE
não prova concorrência temporal: relações, transições e identidade abstrata
pertencem ao Alloy; a transformação textual determinística pertence ao Rust.

\subsection*{Identidade Resumo x Bem}
\texttt{Resumo\_Bem\_Patrimonial} é a identidade catalográfica compartilhável;
\texttt{Bem\_Patrimonial} é a unidade física individual. Editar ou reclassificar
uma unidade não renomeia implicitamente o resumo compartilhado das demais: uma
alteração individual de nome é reclassificação para resumo existente ou criação
explícita de outro resumo. Há testemunha de dois bens com o mesmo resumo e de
reclassificação individual que preserva o resumo do outro bem.

\subsection*{Plaqueta canônica e unicidade permanente}
A plaqueta é a identidade externa da unidade. A forma canônica é
\[ N(s) = s.\mathrm{trim}().\mathrm{toUpperCase}(), \]
não vazia e com no máximo 30
caracteres; zeros iniciais, espaços internos e pontuação são preservados. O
Alloy representa a identidade já canonicalizada pela relação abstrata
\texttt{Plaqueta}: uma plaqueta canônica não identifica dois bens distintos.
\texttt{Chaves\_Unicas} é modelada como reserva permanente (\texttt{reservas}),
distinta de lock: a baixa não remove o bem nem libera a plaqueta para reuso.

\subsection*{Máquina de estado e conservação ortogonal}
O ciclo V1 fechado é \texttt{Ativo -> Inservivel -> Ja\_dado\_baixa}. São
impossíveis o salto \texttt{Ativo -> Ja\_dado\_baixa}, a reversão
\texttt{Inservivel -> Ativo} e qualquer mutação de negócio posterior ao estado
terminal. \texttt{Ja\_dado\_baixa} é terminal para fatos da unidade, mas o bem
permanece existente para retenção, leitura e auditoria. A conservação
(\texttt{BOM}/\texttt{REGULAR}/\texttt{RUIM}) é ortogonal ao status: há
testemunha de bem \texttt{Inservivel} com conservação \texttt{BOM}, impedindo
que o formalizador introduza uma equivalência inexistente.

\subsection*{Versão}
\texttt{versao} começa em 1 e incrementa exatamente uma vez em cada mutação
canônica da unidade (edição aprovada, reclassificação, mudança de local, foto,
conservação, status e baixa). Conflito de versão não altera o bem. Fan-out
derivado de \texttt{nome\_equipamento} ou das projeções de \texttt{Local} não
incrementa a versão e não cria evento de negócio, evitando conflito otimista
falso.

\subsection*{Requisições e locks}
Há no máximo uma edição pendente por bem (\texttt{bem\_edicao\_\{idBem\}}) e
uma adição pendente por plaqueta canônica
(\texttt{bem\_adicao\_\{numeroNormalizado\}}). Todo lock pertence
inequivocamente a uma requisição e carrega tipo e chave de recurso; uma
requisição não remove lock alheio, lock incompatível impede mutação e lock
ausente no momento da resposta é falha de integridade que não permite aprovar
nem rejeitar como íntegra. Nos desfechos terminais, o lock da própria
requisição deixa de existir.

\subsection*{Histórico e baixa}
Cadastro, edição canônica aprovada e baixa válida produzem exatamente um evento
histórico (\texttt{cadastro}, \texttt{edicao}, \texttt{baixa}); retry M7 não
cria fato novo e o fan-out derivado não cria histórico de negócio. A baixa é
rito próprio: exige bem \texttt{Inservivel}, autorização M9 e comprovante
válido; sem comprovante válido a baixa não ocorre. O resultado é
\texttt{Ja\_dado\_baixa}, versão incrementada uma vez, comprovante vinculado,
evento \texttt{baixa}, plaqueta e bem preservados.

\subsection*{Composição M7 e M9}
A composição com M7 é abstrata: um retry do mesmo comando não produz bem,
evento, baixa nem notificação extras. A composição com M9 assume a precondição
\texttt{podeExecutar}/\texttt{podeCommitar}: sem autorização válida no estado do
commit não há mutação patrimonial crítica. Professor não aprova requisição nem
executa baixa; Chefia e Gestor autorizados ainda precisam satisfazer a
precondição de domínio. Autorização é necessária, nunca suficiente.

\subsection*{Responsabilidade das camadas}
CUE valida shapes e limites locais. Alloy valida identidade, unicidade,
máquina de estados, terminalidade, versão, locks, histórico, baixa e composição
abstrata. Rust valida proveniência, hashes, conjunto ordenado do receipt e a
canonicalização determinística $N(s)$ preservando zeros iniciais. O guard
mecânico compara enums, nomes e predicados centrais entre documentação, CUE,
Alloy e Rust.

\begingroup\small\raggedright
% Gerado por lcqui-spec-doc; não editar.
\subsection*{M10\_Patrimonio --- formalização executável M10}
Shapes estruturais de bem, resumo catalográfico, requisições de adição/edição, lock, baixa, evento histórico e alteração. CUE verifica campos, enums, nulabilidade e limites locais; identidade Resumo x Bem, máquina de estados, locks, unicidade permanente, versionamento, terminalidade e composição M7/M9 são verificados em Alloy; canonicalização de plaqueta e proveniência por Rust. CUE não prova concorrência temporal.
\begin{description}
\item[id\_resumo\_bem\_patrimonial] Tipo: string. Obrigatório: sim. Aceita nulo: não.
SQL: TEXT.
FK do modelo catalográfico compartilhado.
\item[numero\_patrimonio] Tipo: string. Obrigatório: sim. Aceita nulo: não.
SQL: TEXT.
Plaqueta canônica trim().toUpperCase(), não vazia, no máximo 30 caracteres.
\item[estado\_conservacao] Tipo: enum. Obrigatório: sim. Aceita nulo: não.
SQL: ENUM.
Valores: BOM, REGULAR, RUIM.
\item[id\_local] Tipo: string. Obrigatório: sim. Aceita nulo: não.
SQL: TEXT.
Autoridade de localização; predio/andar/sala são projeções derivadas.
\item[photo\_url] Tipo: string. Obrigatório: sim. Aceita nulo: não.
SQL: TEXT.
Foto obrigatória da unidade física; validação binária pertence ao backend.
\item[documento\_dado\_baixa\_pdf\_url] Tipo: string. Obrigatório: sim. Aceita nulo: sim.
SQL: TEXT.
Nulo antes da baixa; vinculado e imutável no estado terminal.
\item[nome\_responsavel\_sei] Tipo: string. Obrigatório: sim. Aceita nulo: não.
SQL: TEXT.
Responsável institucional; no máximo 150 caracteres.
\item[status] Tipo: enum. Obrigatório: sim. Aceita nulo: não.
SQL: ENUM.
Valores: Ativo, Inservivel, Ja\_dado\_baixa.
\item[descricao\_complementar] Tipo: string. Obrigatório: sim. Aceita nulo: sim.
SQL: TEXT.
Observação da unidade física; no máximo 200 caracteres.
\item[versao] Tipo: int. Obrigatório: sim. Aceita nulo: não.
SQL: INTEGER.
Inicia em 1; incrementa somente em fato canônico, nunca em fan-out derivado.
\item[numero\_patrimonio\_proposto] Tipo: string. Obrigatório: sim. Aceita nulo: não.
SQL: TEXT.
Plaqueta bruta informada na requisição de adição.
\item[numero\_patrimonio\_normalizado] Tipo: string. Obrigatório: sim. Aceita nulo: não.
SQL: TEXT.
N(proposto) = trim().toUpperCase(); única chave de lock e Chaves\_Unicas.
\item[versao\_bem\_origem] Tipo: int. Obrigatório: sim. Aceita nulo: não.
SQL: INTEGER.
Versão observada para controle otimista de edição.
\item[id\_requisicao] Tipo: string. Obrigatório: sim. Aceita nulo: não.
SQL: TEXT.
Requisição proprietária do lock; precisa coincidir antes da liberação.
\item[tipo\_lock] Tipo: enum. Obrigatório: sim. Aceita nulo: não.
SQL: ENUM.
Valores: EDICAO, ADICAO.
\item[chave\_recurso] Tipo: string. Obrigatório: sim. Aceita nulo: não.
SQL: TEXT.
ID do bem (edição) ou plaqueta normalizada (adição).
\item[tipo\_historico] Tipo: enum. Obrigatório: sim. Aceita nulo: não.
SQL: ENUM.
Valores: cadastro, edicao, baixa.
\end{description}

% Gerado por lcqui-spec-doc; não editar.
\subsection*{Evidência formal M10 --- patrimônio}
CUE verifica shapes de bem, resumo catalográfico, requisições de adição/edição, lock, baixa, evento histórico e alteração, com enums fechados, nulabilidade e limites de comprimento. Alloy verifica a identidade Resumo x Bem, a plaqueta canônica única e não reutilizável, Chaves\_Unicas como reserva permanente distinta de lock, a máquina V1 Ativo -> Inservivel -> Ja\_dado\_baixa, a terminalidade, a ortogonalidade da conservação, o versionamento canônico (fan-out derivado não incrementa), os conflitos de versão/unicidade com liberação do próprio lock, a presença/ownership do lock, o histórico cadastro/edição/baixa e a composição abstrata com M7 (retry não duplica) e M9 (autorização necessária, domínio não dispensado). Rust valida proveniência, ordem exata do receipt e a canonicalização determinística N(s)=trim().toUpperCase() preservando zeros iniciais. A evidência não certifica Firebase, backend, Rules, Storage, assinatura real de PDF, Chaves\_Unicas/backfill, índices, migração de histórico legado nem concorrência sob carga.
\begin{description}
\item[M10-INV-001] PlaquetaUnicaPorBem: UNSAT.\newline Escopo: \texttt{check PlaquetaUnicaPorBem for 6}.
\item[M10-INV-002] ChaveUnicaImpedeNovoBem: UNSAT.\newline Escopo: \texttt{check ChaveUnicaImpedeNovoBem for 6}.
\item[M10-INV-003] BaixaNaoLiberaChave: UNSAT.\newline Escopo: \texttt{check BaixaNaoLiberaChave for 6}.
\item[M10-INV-004] PlaquetaNaoReutilizadaAposBaixa: UNSAT.\newline Escopo: \texttt{check PlaquetaNaoReutilizadaAposBaixa for 8}.
\item[M10-INV-005] ReclassificacaoNaoRenomeiaResumoCompartilhado: UNSAT.\newline Escopo: \texttt{check ReclassificacaoNaoRenomeiaResumoCompartilhado for 8}.
\item[M10-INV-006] CadastroComecaAtivoVersao1: UNSAT.\newline Escopo: \texttt{check CadastroComecaAtivoVersao1 for 8}.
\item[M10-INV-007] SemSaltoAtivoParaBaixa: UNSAT.\newline Escopo: \texttt{check SemSaltoAtivoParaBaixa for 6}.
\item[M10-INV-008] BaixaExigeComprovante: UNSAT.\newline Escopo: \texttt{check BaixaExigeComprovante for 6}.
\item[M10-INV-009] BaixaTerminal: UNSAT.\newline Escopo: \texttt{check BaixaTerminal for 6}.
\item[M10-INV-010] BaixaPreservaBem: UNSAT.\newline Escopo: \texttt{check BaixaPreservaBem for 6}.
\item[M10-INV-011] MutacaoCanonicaIncrementaUmaVez: UNSAT.\newline Escopo: \texttt{check MutacaoCanonicaIncrementaUmaVez for 6}.
\item[M10-INV-012] FanOutNaoIncrementaVersao: UNSAT.\newline Escopo: \texttt{check FanOutNaoIncrementaVersao for 6}.
\item[M10-INV-013] FanOutLocalNaoIncrementaVersao: UNSAT.\newline Escopo: \texttt{check FanOutLocalNaoIncrementaVersao for 6}.
\item[M10-INV-014] FanOutNaoCriaEvento: UNSAT.\newline Escopo: \texttt{check FanOutNaoCriaEvento for 6}.
\item[M10-INV-015] ConflitoVersaoNaoAlteraBem: UNSAT.\newline Escopo: \texttt{check ConflitoVersaoNaoAlteraBem for 6}.
\item[M10-INV-016] RetryNaoDuplicaFato: UNSAT.\newline Escopo: \texttt{check RetryNaoDuplicaFato for 6}.
\item[M10-INV-017] UmaEdicaoPendentePorBem: UNSAT.\newline Escopo: \texttt{check UmaEdicaoPendentePorBem for 6}.
\item[M10-INV-018] UmaAdicaoPendentePorPlaqueta: UNSAT.\newline Escopo: \texttt{check UmaAdicaoPendentePorPlaqueta for 6}.
\item[M10-INV-019] LockTemRequerente: UNSAT.\newline Escopo: \texttt{check LockTemRequerente for 6}.
\item[M10-INV-020] LockAlheioNaoRemovido: UNSAT.\newline Escopo: \texttt{check LockAlheioNaoRemovido for 6}.
\item[M10-INV-021] LockAusenteImpedeResposta: UNSAT.\newline Escopo: \texttt{check LockAusenteImpedeResposta for 6}.
\item[M10-INV-022] CadastroGeraUmEvento: UNSAT.\newline Escopo: \texttt{check CadastroGeraUmEvento for 8}.
\item[M10-INV-023] EdicaoGeraUmEvento: UNSAT.\newline Escopo: \texttt{check EdicaoGeraUmEvento for 6}.
\item[M10-INV-024] BaixaGeraUmEvento: UNSAT.\newline Escopo: \texttt{check BaixaGeraUmEvento for 6}.
\item[M10-INV-025] M9InvalidoImpedeCommit: UNSAT.\newline Escopo: \texttt{check M9InvalidoImpedeCommit for 8}.
\item[M10-INV-026] ProfessorNaoAprova: UNSAT.\newline Escopo: \texttt{check ProfessorNaoAprova for 8}.
\item[M10-INV-027] ProfessorNaoBaixa: UNSAT.\newline Escopo: \texttt{check ProfessorNaoBaixa for 6}.
\item[M10-INV-028] AutorizacaoNaoDispensaDominio: UNSAT.\newline Escopo: \texttt{check AutorizacaoNaoDispensaDominio for 6}.
\item[M10-INV-029] TransicoesPreservamCoerencia: UNSAT.\newline Escopo: \texttt{check TransicoesPreservamCoerencia for 6}.
\item[M10-WIT-030] WitnessDoisBensMesmoResumo: SAT.\newline Escopo: \texttt{run WitnessDoisBensMesmoResumo for 6}.
\item[M10-WIT-031] WitnessReclassificacaoIndividual: SAT.\newline Escopo: \texttt{run WitnessReclassificacaoIndividual for 8}.
\item[M10-WIT-032] WitnessBemAtivo: SAT.\newline Escopo: \texttt{run WitnessBemAtivo for 4}.
\item[M10-WIT-033] WitnessTransicaoInservivel: SAT.\newline Escopo: \texttt{run WitnessTransicaoInservivel for 6}.
\item[M10-WIT-034] WitnessTransicaoBaixa: SAT.\newline Escopo: \texttt{run WitnessTransicaoBaixa for 6}.
\item[M10-WIT-035] WitnessCadastroAprovado: SAT.\newline Escopo: \texttt{run WitnessCadastroAprovado for 6}.
\item[M10-WIT-036] WitnessEdicaoAprovada: SAT.\newline Escopo: \texttt{run WitnessEdicaoAprovada for 6}.
\item[M10-WIT-037] WitnessEdicaoRejeitada: SAT.\newline Escopo: \texttt{run WitnessEdicaoRejeitada for 6}.
\item[M10-WIT-038] WitnessConflitoVersao: SAT.\newline Escopo: \texttt{run WitnessConflitoVersao for 6}.
\item[M10-WIT-039] WitnessConflitoUnicidade: SAT.\newline Escopo: \texttt{run WitnessConflitoUnicidade for 6}.
\item[M10-WIT-040] WitnessDuasPlaquetasDistintas: SAT.\newline Escopo: \texttt{run WitnessDuasPlaquetasDistintas for 6}.
\item[M10-WIT-041] WitnessLockIntegro: SAT.\newline Escopo: \texttt{run WitnessLockIntegro for 6}.
\item[M10-WIT-042] WitnessFanOutSemVersao: SAT.\newline Escopo: \texttt{run WitnessFanOutSemVersao for 6}.
\item[M10-WIT-043] WitnessRetrySemDuplicacao: SAT.\newline Escopo: \texttt{run WitnessRetrySemDuplicacao for 6}.
\item[M10-WIT-044] WitnessBaixaComHistorico: SAT.\newline Escopo: \texttt{run WitnessBaixaComHistorico for 6}.
\item[M10-WIT-045] WitnessBemBaixadoExiste: SAT.\newline Escopo: \texttt{run WitnessBemBaixadoExiste for 6}.
\item[M10-WIT-046] WitnessConservacaoOrtogonal: SAT.\newline Escopo: \texttt{run WitnessConservacaoOrtogonal for 4}.
\end{description}
UNSAT para check indica ausência de contraexemplo no escopo declarado; SAT para run indica testemunha encontrada. A evidência não certifica a implementação Firebase/backend.

\par\endgroup

\subsection*{Limites}
Esta evidência prova somente o modelo formal no escopo declarado e \textbf{não}
certifica \texttt{functions/}\allowbreak\texttt{src/}\allowbreak\texttt{patrimonio.ts}, o frontend, \texttt{firestore.rules},
\texttt{storage.rules}, Firebase, Admin SDK, IAM, App Check, transações reais
sob carga, o Storage real, a assinatura real dos arquivos, o backfill de
\texttt{Chaves\_Unicas}, índices, a migração do histórico legado nem a
infraestrutura em produção. A prova de consistência do modelo não é homologação
da implementação real.
\fussy
